Temporal point processes are used both to \emph{measure} market mechanisms, for example how fast
prices absorb news and how much of the reaction is the market responding to itself, and to
\emph{predict} the next event. We show that one algebraic object, the phase-type (Erlang-chain)
state space, serves both purposes with exact likelihoods. For measurement, we build a multiscale
phase-type Hawkes process whose parameters all enter linearly. Its log-likelihood is concave, and
we compute the global maximum together with a Fenchel duality-gap certificate. A Markov embedding
gives the response to an exogenous shock, echoes included, in closed form, and this response is
stable if and only if the spectral radius of the branching matrix is below one. Applied to six
markets around 670 scheduled macroeconomic releases (2022--2026) with a built-in placebo test, the
model shows that the direct reaction to CPI, payrolls, PPI and the FOMC statement is absorbed in
about a second, while the market's own echoes stretch half of the response to
\tEchoLo--\tEchoHi{}\,s. For prediction, we introduce EPT-X, a neural point process that adds a
history-dependent mixture-of-Erlang renewal channel to a linear Hawkes backbone. Its intensity
and compensator remain in closed form, so it is trained and evaluated on the exact likelihood.
EPT-X has the best test likelihood on all five LOBSTER order-book streams against six classical
estimators and seven neural baselines, improves the best published result on Retweet by more
than 0.6~nats per event, is within one standard deviation of it on Taobao and StackOverflow, and
trails the state-space model S2P2 on Taxi and Amazon. We also find that a widely used
IntensityFree implementation clamps short inter-event times before scoring, which overstates its
likelihood on streams with many short gaps.
